\documentclass[12pt]{article}

% Packages
\usepackage{amsmath, amssymb}
\usepackage{geometry}
\usepackage{graphicx}
\graphicspath{{../../images/}}

\usepackage[colorlinks = true,
            linkcolor = blue,
            urlcolor  = blue,
            citecolor = blue,
            anchorcolor = blue]{hyperref}
\usepackage{enumitem}
\usepackage[backend=biber, style=numeric]{biblatex}
\addbibresource{sba_classification.bib}
\usepackage[skins]{tcolorbox}
\usepackage[T1]{fontenc}
\usepackage[ttdefault=true]{AnonymousPro}
\renewcommand*\familydefault{\ttdefault} %% Only if the base font of the document is to be typewriter style
% Page setup

\geometry{margin=1in}





\begin{document}

\begin{tcolorbox}[halign=flush center]
         \textbf{Representation of Good Borrowers using Graph Embeddings}
   
 \end{tcolorbox}

\section{Introduction}
Developing a representations of data that is part of a graph so that it can be used in a machine learning model is an important problem. In this example, we consider the problem of developing a representation for good borrowers in the SBA 7a loans dataset. This representation can be used for tasks such as developing a taxonomy of good borrowers. Understanding the heterogeneity of good borrowers can inform the design of new loan products. It can also help identify good borrowers that are similar to bad borrowers. This can help identify potential risks in the loan portfolio. In the previous exercise, graphs were used to develop features for borrowers. In this example, we represent good borrowers as a graph and then develop a representation for this graph using a spectral embedding technique. The representation can then be used in a machine learning model.
We start with a set of relations and then develop a graph. Spectral embedding is then used to develop a representation for the graph.

.

\printbibliography

\end{document}